\documentclass[12pt,a4paper]{article}
\usepackage[utf8]{inputenc}
\usepackage[T1]{fontenc}
\usepackage[italian]{babel}
\usepackage{libertinus}
\usepackage{amssymb}
\usepackage[a4paper,top=25mm,bottom=25mm,left=22mm,right=22mm,headheight=15pt]{geometry}
\usepackage{graphicx}
\graphicspath{{../../assets/}{../../Assets/}{../assets/}{../Assets/}{assets/}{Assets/}}
\usepackage[table]{xcolor}
\usepackage{tabularx,array}
\usepackage{titlesec}
\usepackage{titletoc}
\usepackage{microtype}
\usepackage{fancyhdr}
\usepackage{enumitem}
\usepackage[hidelinks]{hyperref}

% Numero del verbale: da impostare una sola volta (usato in intestazione, copertina e ID)
\newcommand{\vnum}{6}

% Palette colori del gruppo Bending Forks
\definecolor{bfblue}{HTML}{1B3F9E}
\definecolor{bfrow}{HTML}{F2F5FC}
\definecolor{bfline}{HTML}{C9D1E6}
\definecolor{bfhead}{HTML}{2E6BE6}

\newcommand{\tightlist}{\setlength{\itemsep}{1pt}\setlength{\parskip}{0pt}}
\setlist{topsep=2pt}
\setlist[itemize]{leftmargin=1.3em,label=\textbullet}
\setlist[enumerate]{leftmargin=1.6em}
\setlength{\parindent}{0pt}
\sloppy
\setlength{\parskip}{5pt}
\renewcommand{\arraystretch}{1.45}
\setcounter{secnumdepth}{3}
\setcounter{tocdepth}{3}

% Formattazione titoli sezioni
\titleformat{\section}{\LARGE\bfseries\color{bfblue}}{\thesection.}{0.3em}{}[{\color{bfblue}\vspace{-2pt}\titlerule[0.9pt]}]
\titleformat{\subsection}{\Large\bfseries\color{bfblue}}{\thesubsection.}{0.3em}{}
\titleformat{\subsubsection}{\normalsize\bfseries}{\thesubsubsection.}{0.3em}{}
\titlespacing*{\section}{0pt}{14pt}{6pt}
\titlespacing*{\subsection}{0pt}{10pt}{4pt}
\titlespacing*{\subsubsection}{0pt}{10pt}{3pt}

% Indice cliccabile con puntini/trattini
\contentsmargin{2.2em}
\newcommand{\bfleader}{\titlerule*[0.32em]{\textendash}}
\titlecontents{section}[0pt]{\addvspace{5pt}\large\bfseries}{\thecontentslabel.\ }{}{\hspace{0.4em}\bfleader\contentspage}
\titlecontents{subsection}[1.2em]{\addvspace{0pt}}{\thecontentslabel.\ }{}{\hspace{0.4em}\bfleader\contentspage}
\titlecontents{subsubsection}[2.4em]{\addvspace{0pt}}{\thecontentslabel.\ }{}{\hspace{0.4em}\bfleader\contentspage}

% Intestazioni e piè di pagina
\pagestyle{fancy}
\fancyhf{}
\fancyhead[L]{\footnotesize\sffamily\bfseries\color{bfhead}\textls[120]{BENDING FORKS}}
\fancyhead[R]{\footnotesize\sffamily Verbale Interno N. \vnum}
\fancyfoot[C]{\small Pagina \thepage}
\renewcommand{\headrulewidth}{0.9pt}

% Colorazione alternata per le tabelle
\newcommand{\bftable}{\rowcolors{2}{bfrow}{white}}
\usepackage{etoolbox}
\AtBeginEnvironment{tabularx}{\bftable}

\begin{document}

% ---------- Copertina ----------
\begin{titlepage}
\centering
\vspace*{4mm}
\includegraphics[width=105mm]{logo-bending-forks-trasparente.png}\\[4mm]
\underline{bending.forks.swe@gmail.com}

\vspace{42mm}
{\fontsize{24}{28}\selectfont\bfseries Verbale Interno N. \vnum\par}
\vspace{4mm}
{\Large\color{bfblue} Riunione del 4 novembre 2026\par}

\vfill
\begin{tabular}[t]{r|l}
\textbf{Responsabile} & Nicholas Salmaso \\
\textbf{Segretario} & Filippo Rampazzo \\
\textbf{Moderatore} & Filippo Rampazzo \\
\textbf{Verificatore} & Matteo Wu \\
\end{tabular}
\hspace{14mm}
\begin{tabular}[t]{r|l}
\textbf{Versione} & 1.0 \\
\textbf{Tipo} & Verbale Interno \\
\textbf{Destinatari} & \emph{Gruppo Bending Forks} \\
 & \emph{Prof. Tullio Vardanega} \\
 & \emph{Prof. Riccardo Cardin} \\
\end{tabular}

\vspace{12mm}
{\large\emph{\textbf{Abstract dei contenuti:}}}\\[1mm]
\begin{minipage}{150mm}\centering Sintesi e decisioni operative emerse durante la riunione interna del team del giorno 4 novembre 2026, riguardante la pianificazione delle attività e l'avanzamento del progetto.\end{minipage}
\vspace*{6mm}
\end{titlepage}

% ---------- Indice ----------
\vspace*{2mm}
{\fontsize{22}{26}\selectfont\bfseries Indice\par}
\vspace{6mm}
\makeatletter\@starttoc{toc}\makeatother
\newpage

% ---------- Registro delle versioni ----------
\hypertarget{registro-delle-versioni}{%
\section{Registro delle versioni}}

\begin{center}
\begin{tabularx}{\linewidth}{
  >{\hsize=0.55\hsize\raggedright\arraybackslash}X
  >{\hsize=0.75\hsize\raggedright\arraybackslash}X
  >{\hsize=0.85\hsize\raggedright\arraybackslash}X
  >{\hsize=0.85\hsize\raggedright\arraybackslash}X
  >{\hsize=2.00\hsize\raggedright\arraybackslash}X
}
\rowcolor{bfblue}\textcolor{white}{\textbf{Versione}} & 
\textcolor{white}{\textbf{Data}} & 
\textcolor{white}{\textbf{Autore}} & 
\textcolor{white}{\textbf{Verificatore}} & 
\textcolor{white}{\textbf{Descrizione}} \\
1.0 & 06/11/2026 & Filippo Rampazzo & Matteo Wu & Creazione e redazione iniziale del documento \\ \arrayrulecolor{bfline}\hline
\end{tabularx}
\end{center}

% ---------- Informazioni generali ----------
\hypertarget{informazioni-generali}{%
\section{Informazioni generali}}

\hypertarget{durata-e-luogo}{%
\subsection{Durata e luogo}}

\begin{itemize}
\tightlist
\item \textbf{Data:} 4 novembre 2026
\item \textbf{Inizio:} 15:00
\item \textbf{Fine:} 17:30
\item \textbf{Luogo / modalità:} in remoto, tramite Discord
\end{itemize}

\hypertarget{partecipanti}{%
\subsection{Partecipanti}}

\begin{center}
\begin{tabularx}{\linewidth}{>{\raggedright\arraybackslash}X>{\centering\arraybackslash}X>{\centering\arraybackslash}p{2.2cm}>{\centering\arraybackslash}p{2.2cm}}
\rowcolor{bfblue}\textcolor{white}{\textbf{Nome e Cognome}} & \textcolor{white}{\textbf{Ruolo Progetto}} & \textcolor{white}{\textbf{Presente}} & \textcolor{white}{\textbf{Assente}} \\
Marco Beltrame & Amministratore & $\blacksquare$ & $\square$ \\ \arrayrulecolor{bfline}\hline
Lorenzo Milanese & Amministratore & $\blacksquare$ & $\square$ \\ \arrayrulecolor{bfline}\hline
Filippo Rampazzo & Analista & $\blacksquare$ & $\square$ \\ \arrayrulecolor{bfline}\hline
Nicholas Salmaso & Responsabile & $\blacksquare$ & $\square$ \\ \arrayrulecolor{bfline}\hline
Luca Stangherlin & Analista & $\square$ & $\blacksquare$ \\ \arrayrulecolor{bfline}\hline
Aldo Toffoli & Analista & $\blacksquare$ & $\square$ \\ \arrayrulecolor{bfline}\hline
Matteo Wu & Verificatore & $\blacksquare$ & $\square$ \\ \arrayrulecolor{bfline}\hline
\end{tabularx}
\end{center}

% ---------- Ordine del giorno ----------
\hypertarget{ordine-del-giorno}{%
\section{Ordine del giorno}}

\begin{enumerate}
\tightlist
\item \hyperlink{odg-1}{Comunicazione alla proponente dell'aggiudicazione dell'appalto};
\item \hyperlink{odg-2}{Pianificazione della stesura dei prossimi documenti necessari};
\item \hyperlink{odg-3}{Suddivisione dei ruoli di progetto rendicontati};
\item \hyperlink{varie-ed-eventuali}{Varie ed eventuali}.
\end{enumerate}

% ---------- Discussione e decisioni ----------

\newpage
\hypertarget{discussione-e-decisioni}{%
\section{Discussione e decisioni}}

\hypertarget{odg-1}{%
\subsection{Comunicazione alla proponente dell'aggiudicazione dell'appalto}}

\textbf{Discussione:} \\
Il gruppo ha comunicato alla proponente l'aggiudicazione dell'appalto.
Rampazzo e Toffoli propongono di fissare una riunione con la proponente, in modo da discutere i dettagli del progetto e le prossime fasi operative.
Il gruppo concorda di presentarsi alla riunione con un abbozzo di Analisi dei Requisiti, in modo da poter porre domande specifiche.
I partecipanti concordano di non fissare la riunione troppo in avanti per evitare di rimanere bloccati senza riscontro da parte della proponente, ma neanche troppo presto per via del considerevole carico di studio richiesto da una prova parziale il 9 novembre.
Il gruppo si accorda su venerdì 13 novembre 2026 come data che soddisfa le suddette esigenze.

\textbf{Decisioni prese:}
\begin{itemize}
\tightlist
\item[$\checkmark$] Il gruppo ha inviato una email alla proponente comunicando l'aggiudicazione dell'appalto e richiedendo una riunione per venerdì 13 novembre 2026;
\item[$\checkmark$] Il gruppo fissa la riunione con la proponente come scadenza per l'abbozzo dell'Analisi dei Requisiti.
\end{itemize}

\textbf{Motivazione:} \\
La data del 13 novembre consente di stilare l'abbozzo dei documenti necessari, lasciando tempo ai membri del gruppo per concentrarsi sull'imminente prova parziale.

\textbf{Responsabili attuazione:} fissati all'\hyperlink{odg-2}{OdG 2} \\
\textbf{Scadenza:} 13 novembre 2026 (solo bozza Analisi dei Requisiti)

\hypertarget{odg-2}{%
\subsection{Pianificazione della stesura dei prossimi documenti necessari}}

\textbf{Discussione:} \\
Dopo una discussione collettiva sulla loro struttura e contenuti generali, il gruppo individua i seguenti tre documenti da redigere: Analisi dei Requisiti, Norme di Progetto e Glossario.
Toffoli propone di redigere collettivamente in serie tutti e tre i documenti;
Beltrame e Rampazzo propongono di redigere i documenti in parallelo, dividendosi in sotto-gruppi.
Il gruppo concorda sulla divisione in sotto-gruppi, uno per documento, così composti:
\begin{itemize}
\tightlist
\item \textbf{Analisi dei Requisiti:} Rampazzo, Stangherlin, Toffoli;
\item \textbf{Norme di Progetto:} Beltrame, Milanese;
\item \textbf{Glossario:} Salmaso, Wu.
\end{itemize}

\textbf{Decisione presa:}
\begin{itemize}
\tightlist
\item[$\checkmark$] Suddivisione in sotto-gruppi paralleli per la redazione dei documenti;
\end{itemize}

\textbf{Motivazione:} \\
Evitare che ci siano membri del gruppo senza incarichi assegnati.

\textbf{Scadenza:} 13 novembre 2026 (solo bozza Analisi dei Requisiti)

\hypertarget{odg-3}{%
\subsection{Suddivisione dei ruoli di progetto rendicontati}}

\textbf{Discussione:} \\
Vista la necessità di suddividersi il lavoro secondo quanto indicato nel preventivo, il gruppo si spartisce i ruoli su base volontaria, secondo quanto indicato nella tabella 4.3.1.
Essendo ancora in fase di analisi, si esclude di assegnare il ruolo di programmatore e di progettista.
I ruoli vengono assegnati coerentemente alle mansioni approvate all'\hyperlink{odg-2}{OdG 2}.

\hypertarget{tab-4.3.1}{
\begin{center}
\begin{tabularx}{\linewidth}{>{\raggedright\arraybackslash}X>{\raggedright\arraybackslash}X}
\rowcolor{bfblue}\textcolor{white}{\textbf{Nome e Cognome}} & \textcolor{white}{\textbf{Ruolo Progetto}}\\
Marco Beltrame & Amministratore \\ \arrayrulecolor{bfline}\hline
Lorenzo Milanese & Amministratore \\ \arrayrulecolor{bfline}\hline
Filippo Rampazzo & Analista \\ \arrayrulecolor{bfline}\hline
Nicholas Salmaso & Responsabile \\ \arrayrulecolor{bfline}\hline
Luca Stangherlin & Analista \\ \arrayrulecolor{bfline}\hline
Aldo Toffoli & Analista \\ \arrayrulecolor{bfline}\hline
Matteo Wu & Verificatore \\ \arrayrulecolor{bfline}\hline
\end{tabularx}
\textbf{Tabella 4.3.1 - Suddivisione dei ruoli fino al 16 novembre 2026}
\end{center}
}

Il gruppo approva all'unanimità la suddivisione dei ruoli, attendendo eventuale conferma di Stangherlin, che non ha potuto partecipare alla riunione.
I ruoli sono fissati fino al 16 novembre 2026, in modo da coprire il diario di bordo del 9 novembre e la riunione con la proponente;
a partire dal 16 novembre, il gruppo stabilisce di dare ai ruoli durata bisettimanale.

\textbf{Decisioni prese:}
\begin{itemize}
\tightlist
\item[$\checkmark$] Suddivisione dei ruoli come indicato alla \hyperlink{tab-4.3.1}{Tabella 4.3.1};
\item[$\checkmark$] I ruoli entrano subito in vigore e lo rimangono fino al 16 novembre compreso; successivamente i ruoli saranno riassegnati ogni due settimane;
\end{itemize}

\hypertarget{varie-ed-eventuali}{%
\subsection{Varie ed eventuali}}
Il gruppo prende atto dei feedback del committente sulla documentazione finora prodotta per il progetto e discute sulla loro implementazione. \\
Il gruppo decide di fissare al 12 novembre 2026 la prossima riunione, per discutere del lavoro svolto in vista della riunione con la proponente.
\newpage

% ---------- Azioni da intraprendere (TODO) ----------
\hypertarget{azioni-da-intraprendere}{
\section{Azioni da intraprendere (TODO)}}

\begin{center}
\begin{tabularx}{\linewidth}{>{\hsize=2.0\hsize\raggedright\arraybackslash}X >{\hsize=1.0\hsize\raggedright\arraybackslash}X >{\hsize=0.5\hsize\centering\arraybackslash}X >{\hsize=0.5\hsize\centering\arraybackslash}X}
\rowcolor{bfblue}\textcolor{white}{\textbf{Descrizione}} & \textcolor{white}{\textbf{Responsabile}} & \textcolor{white}{\textbf{Scadenza}} & \textcolor{white}{\textbf{Priorità}} \\
Redigere il verbale della riunione del 4 novembre & Rampazzo & 2026-11-12 & Bassa \\ \arrayrulecolor{bfline}\hline
Completare l'Analisi dei Requisiti v0.1 & Rampazzo, Stangherlin, Toffoli & 2026-11-13 & Alta \\ \arrayrulecolor{bfline}\hline
Completare le Norme di Progetto v0.1 & Beltrame, Milanese & 2026-11-16 & Alta \\ \arrayrulecolor{bfline}\hline
Redigere il Glossario & Salmaso, Wu & 2026-11-16 & Media \\ \arrayrulecolor{bfline}\hline
\end{tabularx}
\end{center}

% ---------- Prossima riunione ----------
\hypertarget{prossima-riunione}{%
\section{Pianificazione prossima riunione}}

\begin{itemize}
\tightlist
\item \textbf{Data proposta:} 12 novembre 2026
\item \textbf{Ora prevista:} 15:00
\item \textbf{Luogo / modalità:} in remoto, tramite Discord 
\item \textbf{OdG preliminare:} Allineamento preliminare in vista della riunione con la proponente programmata per il 13 novembre 2026
\end{itemize}

\vspace{35pt}
\noindent
\begin{minipage}[t]{0.45\textwidth}
\textbf{Firma Scriba:}\\[18pt]
\makebox[0pt][l]{\raisebox{0mm}{\includegraphics[width=40mm]{firma_filippo.jpeg}}}\rule{55mm}{0.5pt}\\[2pt]
\small Filippo Rampazzo
\end{minipage}
\hfill
\begin{minipage}[t]{0.45\textwidth}
\begin{flushright}
\textbf{Firma Verificatore:}\\[18pt]
\makebox[0pt][l]{\raisebox{0mm}{\includegraphics[width=40mm]{firma_matteo.jpg}}}\rule{55mm}{0.5pt}\\[2pt]
\small Matteo Wu
\end{flushright}
\end{minipage}

\end{document}